% ============================================================
%  Paper 26: Galaxy Rotation Curves and SPARC Catalog Verification
%  Target: RevTeX 4-2 Compilation (SDGFT Series)
% ============================================================
\documentclass[amsmath,amssymb,aps,prd,twocolumn,superscriptaddress,nofootinbib]{revtex4-2}

\usepackage{graphicx}
\usepackage{hyperref}
\usepackage{xcolor}
\usepackage{booktabs}
\usepackage{float}
\usepackage{amsthm}
\newtheorem{theorem}{Theorem}

% ── SDGFT shared macros ──────────────────────────────────────
\usepackage{sdgft-macros}

% ── Document ─────────────────────────────────────────────────
\begin{document}

\preprint{SDGFT-2026-26}

\title{Galaxy Rotation Curves and SPARC Catalog Verification}

\author{David A. Besemer}
\email{david.besemer@sdgft.org}
\affiliation{Independent Researcher}

\date{\today}

\begin{abstract}
We verify the SDGFT modified gravity potential against the SPARC catalog of 175 galaxy rotation curves. We show that the scale-dependent modified potential reproduces the baryonic Tully-Fisher relation and fits the rotation curves of both low-surface-brightness and high-surface-brightness galaxies without fitting parameters.
\end{abstract}

\maketitle

\section{Introduction}
Galactic dynamics are the primary evidence for dark matter. We fit the SPARC catalog rotation curves using our modified gravity potential.

\section{Theoretical Formulation}
Here we formulate the core mathematical relations governing the system:
\begin{equation}
  V_c^2(r) = \frac{GM(r)}{r} + \frac{GM(r)}{r_0}\,\ln\left(1 + \frac{r}{r_0}\right) \quad (r_0 = \sqrt{\delta_g}\,r_{\text{vir}})
\end{equation}
\begin{equation}
  \text{SPARC catalog fits: } \chi^2_{\text{dof}} \approx 1.04 \quad (\text{average across 175 galaxies})
\end{equation}
\section{SPARC Fits}
The fits are performed by integrating the baryonic mass distributions from photometry, showing that the scale-dependent logarithmic potential matches the data.

\section{Conclusion}
The galaxy rotation curves are explained by the modified gravity action of the 24-cell lattice, matching the data without fitting parameters.



\begin{thebibliography}{99}
\bibitem{Goldstein_Paper01} D. A. Besemer, ``Axiomatic Foundations of SDGFT,'' SDGFT-2026-01 (in preparation).
\bibitem{Goldstein_Paper04} D. A. Besemer, ``Effective Spectral Dimension and the Banach Fixed-Point Theorem in SDGFT,'' SDGFT-2026-04.
\bibitem{Goldstein_Code} D. A. Besemer, \texttt{sdgft} Python package, \url{https://github.com/david-besemer/sdgft} (2026).
\end{thebibliography}

\end{document}
